\medskip
































\medskip

\textbf{PARTIE I}

\medskip

\begin{enumerate}
\item On traduit la situation à l'aide d'un arbre pondéré:
\begin{center}
\pstree[treemode=R,nodesepB=3pt,levelsep=2.75cm]{\TR{}}
{\pstree{\TR{$D$~~}\taput{0,05}}
	{\TR{$T$}\taput{0,98}
	\TR{$\overline{T}$}\tbput{0,02}
	}
\pstree{\TR{$\overline{D}$~~}\tbput{0,95}}
	{\TR{$T$}\taput{0,03}
	\TR{$\overline{T}$}\tbput{0,97}
	}	
}
\end{center}

\item \begin{enumerate}
\item 
On a: $P(D \cap T) = P(D) \times P_D(T) = 0,05 \times 0,98 = 0,049$.
\item 
On a de même : $P\left(\overline{D} \cap T\right) = P\left(\overline{D}\right)  \times P_{\overline{D}}(T) = 0,95 \times 0,3 = \np{0,0285}$.

D'après la loi des probabilités totales :

$P(T) = P(D \cap T) + P\left(\overline{D} \cap T\right) = 0,049 + \np{0,0285} = \np{0,0775}$.
\end{enumerate}
\item 


La valeur prédictive positive du test est égale à :

$P_T(D) = \dfrac{P(T \cap D)}{P(T)} = \dfrac{P(D \cap T)}{P(T)} = \dfrac{0,049}{\np{0,0775}} \approx \np{0,6322}$, soit 0,632 au millième près.

Comme $0,632 < 0,95$ on peut en déduire que le test n'est pas efficace.
\end{enumerate}
\medskip

\textbf{PARTIE II}

\medskip







\begin{enumerate}
\item 
Le choix de l'échantillon étant assimilé à un tirage avec remise et avec une probabilité constante de choisir un produit défectueux égale à $ 0,05$, on peut donc dire que la variable aléatoire $X$ suit une loi binomiale de paramètres $n = 20$ et $p  = 0,05$.
\item 


On a: $p(X = 0) = 0,05^0 \times 0,95^{20}$.

Donc la probabilité cherchée est $p(X \geqslant 1) = 1 -p(X =0) = 1 - 0,95^{20} \approx 0,642$, soit 0,64 au centième près.
\item 
On a: $E = n \times p = 20 \times 0,05 = 1$.

Cela signifie que sur un grand nombre de tirages d'échantillons on trouvera 1 pièce défectueuse sur 20 pièces tirées.
\end{enumerate}

\vspace{0,5cm}
